\subsubsection{Evaluierungsergebnisse der Spitzendetektion}

Die quantitative Auswertung der beiden Testgebiete, deren Ergebnisse in Tabelle \ref{tab:evaulierung} gegenübergestellt sind, zeigt deutliche strukturelle Unterschiede in der Genauigkeit der Baumdetektion.

\begin{table}[htbp]
\centering
\caption{Ergebnisse der Einzelbaumdetektion in den Testgebieten Gohliser Straße und Mariannenpark}
\label{tab:evaluierung}
\begin{tabular}{l c c c c c c}
\toprule
\textbf{} & \textbf{TP} & \textbf{FP} & \textbf{FN} & \textbf{Precision} & \textbf{Recall} & \textbf{F1-Score} \\
\midrule
\textbf{Gohliser Straße} & 45 & 0 & 4 & 0,75 & 0,92 & 0,83 \\
\addlinespace
\textbf{Mariannenpark} & 261 & 24 & 141 & 0,82 & 0,65 & 0,73 \\
\bottomrule
\end{tabular}
\end{table}

Der Vergleich liefert eindeutige Erkenntnisse über die Leistung des Algorithmus bei unterschiedlichen Bestandsdichten. Im lichten Bereich der Straßenbäume (Gohliser Straße) wird ein hoher Recall von 0,92 erzielt, was die zuverlässige Erfassung fast aller registrierten Bäume belegt. Demgegenüber steht eine geringere Precision von 0,75, woraus ein F1-Score von 0,83 resultiert. Während die isolierte Stellung der Straßenbäume die Erkennung begünstigt und die Zahl der nicht erfassten Objekte minimiert, reagiert die Suche nach lokalen Maxima sensibel auf strukturelle Unregelmäßigkeiten in Kronendächern. Bei ausladenden Laubbäumen werden sekundäre Äste fälschlicherweise als eigenständige Baumspitzen detektiert, was zu zusätzlichen FP führt und die Precision senkt.

Im dichten Bestand des Mariannenparks sinkt der F1-Score auf 0,73, wobei sich ein entgegengesetztes Verhalten zeigt. Die Precision liegt bei einem hohen Wert von 0,82, der Recall fällt jedoch auf 0,65, wodurch ein erheblicher Teil des tatsächlichen Bestands unerfasst bleibt. Die Ursache hierfür ist das geschlossene Kronendach der Parkanlage. Wenn Baumkronen nahtlos ineinandergreifen, fehlen im CHM die topografischen Vertiefungen, sodass die Suche nach lokalen Maxima primär dominante Oberkronen detektiert. Dies führt zu einer Untersegmentierung, bei der engstehende Baumgruppen als ein einzelnes Objekt erfasst werden. Da astbedingte Sekundärmaxima durch das geschlossene Kronendach gedämpft werden, bleibt die Überdetektion im Park gering.

Diese Abhängigkeit der Genauigkeit von der Bestandsdichte deckt sich mit aktuellen Studien zur urbanen Forstinventur. Cetin und Yastikli (2025) dokumentieren bei einer punktbasierten Baumdetektion mittels Clustering in Istanbul ein analoges Verhalten mit einer Vollständigkeit von 87,8~\% bei freistehenden Bäumen, während der F1-Score in dichten urbanen Gebieten auf Werte zwischen 0,68 und 0,70 sinkt. Diese Verschlechterung wird ebenfalls auf die fehlerhafte Zusammenfassung ineinandergreifender Kronen zurückgeführt. Dies verdeutlicht, dass dichte Kronendächer unabhängig von der gewählten Methodik eine universelle Herausforderung bei der Einzelbaumdetektion darstellen \cite{cetinAutomaticDetectionUrban2025}.

Zusammenfassend lässt sich festhalten, dass die Leistung des rasterbasierten Workflows direkt von der Geometrie des Kronendachs beeinflusst wird. In lichten Strukturen arbeitet die Detektion lokaler Maxima sehr zuverlässig, neigt dort jedoch zu einer Überdetektion. In dichten Parkstrukturen stößt der Ansatz an Grenzen, da die strukturelle Komplexität mehrschichtiger Bestände durch ein planares Rastermodell nicht vollständig abgebildet werden kann.